\fancyhead[L]{\small ACHROMATISCHE DOPPELOBJEKTIVE.}
\SourceLeaf{510}\pdfbookmark[0]{Über achromatische Doppelobjektive}{achromatische_doppelobjective}
\begin{center}
{\Large ÜBER DIE ACHROMATISCHEN DOPPELOBJEKTIVE\par}
\vspace{6mm}
{\large BESONDERS IM HINBLICK AUF\par}
\vspace{4mm}
{\large DIE VOLLSTÄNDIGE AUFHEBUNG DER FARBZERSTREUUNG.\par}
\vspace{7mm}\hrule\vspace{2mm}
Zeitschrift für Astronomie und verwandte Wissenschaften\par
herausgegeben von B.\ von \textsc{Lindenau} und \textsc{Bohnenberger}. Bd.\ IV. Nr.\ XXX. Dezember 1817.
\vspace{2mm}\hrule\vspace{5mm}
\end{center}
Der schöne Aufsatz von Professor \textsc{Bohnenberger} über achromatische Objektive im ersten Band dieser Zeitschrift hat das Verdienst, einen für diese Theorie wichtigen Umstand zuerst zur Sprache gebracht zu haben. Er hat mich dazu veranlasst, einige frühere Untersuchungen wieder aufzunehmen und weiterzuführen. Ihre Ergebnisse werde ich hier mitteilen.

Bisher begnügte man sich bei Doppelobjektiven damit, die Farbzerstreuung für die der Achse unendlich nahen Strahlen und die Abweichung durch die Kugelgestalt für Strahlen mittlerer Brechbarkeit aufzuheben. Für die Randstrahlen kann dabei noch eine geringe Farbzerstreuung bleiben. Bei dieser Anordnung ist die Berechnung eines achromatischen Objektivs eine unbestimmte Aufgabe. Zu jeder Kronglaslinse mit positiver Brennweite lässt sich also, unabhängig vom Verhältnis der Radien ihrer Flächen, eine Flintglaslinse berechnen, die mit ihr zusammen ein im obigen Sinn achromatisches Objektiv ergibt. Soweit ich weiß, haben bisher alle Optiker beide Flächen der Kronglaslinse als konvex angenommen. Für das Verhältnis der beiden Radien haben die Theoretiker jedoch sehr unterschiedliche Werte vorgeschlagen, je nachdem, von welchem Prinzip sie ausgegangen sind. Will man mit \textsc{Euler} die durch die Gestalt bedingte Abweichung bei der Kronglaslinse möglichst klein machen, müssen die Radien ungefähr im Verhältnis
\SourceLeaf{511}1 zu 7 stehen. Sie müssen gleich sein, wenn man wie \textsc{Klügel} in seiner analytischen Dioptrik die kleinstmöglichen Krümmungen haben möchte. Sollen die Brechungen selbst so klein wie möglich werden, wie es sich derselbe Verfasser in einer späteren Abhandlung vornimmt, so müssen sie einander gleich sein, und die Radien müssen ungefähr im Verhältnis 1 zu 3 stehen. Es scheint nicht, dass alle diese verschiedenen Vorschläge hinreichend begründet sind. \textsc{Klügels} Augenmerk galt besonders der Abweichung durch die Kugelgestalt, die sich bekanntlich unmöglich für \textit{alle} Strahlen mathematisch exakt aufheben lässt. Bei \textsc{Eulers} Behandlung dieser Rechnungen ist die Abweichung eigentlich nur für die achsennächsten Strahlen aufgehoben. Die verbleibende Abweichung ist der vierten Potenz des Abstands von der Achse sehr nahe proportional und macht sich daher bei den Randstrahlen am stärksten bemerkbar. Oder man führt die Rechnung wie \textsc{Klügel} so, dass die Abweichung für die Randstrahlen verschwindet. Dann zeigt sie sich wieder bei den Strahlen dazwischen, am stärksten bei denen, deren Abstand von der Achse annähernd $\tfrac{7}{10}$ des Öffnungsradius beträgt. \textsc{Klügels} Absicht bei der Wahl des Verhältnisses der beiden ersten Radien war es, diese unvermeidlich verbleibende Abweichung durch die Gestalt möglichst unschädlich zu machen. Es geht jedoch weder klar genug hervor, dass das gewählte Verhältnis diesen Zweck tatsächlich am besten erfüllt, noch, dass dieser Zweck wichtig genug ist, um ihn vorrangig und allein der Bestimmung dieses Verhältnisses zugrunde zu legen. Hat ein solches Objektiv nämlich noch \textit{andere} Unvollkommenheiten, die erheblich größer sind als die durch die nicht ganz aufhebbare Abweichung aufgrund der Gestalt verursachten, dann ist es offensichtlich wichtiger, jene als diese zu berücksichtigen.

Deshalb wird es vorteilhafter sein, die Freiheit bei der Bestimmung des Verhältnisses der ersten beiden Radien zu nutzen, um die Farbzerstreuung bei den Randstrahlen zu verringern oder aufzuheben. Tatsächlich hat Professor \textsc{Bohnenberger} durch Rechnung gezeigt, dass in dieser Hinsicht das Verhältnis 2 zu 3 dem Verhältnis 1 zu 3 vorzuziehen ist. Beim ersten ist die Farbzerstreuung der Randstrahlen beträchtlich kleiner, ohne dass die verbleibende Abweichung durch die Kugelgestalt erheblich wird. Allerdings bleibt auch bei Professor \textsc{Bohnenbergers} Anordnung eine Farbzerstreuung der Randstrahlen zurück,
\SourceLeaf{512}die man sehr gerne weiter verringern oder ganz beseitigen würde. Wie Professor \textsc{Bohnenberger} auf S.\ 392 schreibt, waren seine Versuche zu diesem Zweck erfolglos. Sie könnten die Vermutung nahezulegen scheinen, dies sei unmöglich. Das hat mich zu einer eigenen Untersuchung veranlasst, aus der sich das Gegenteil ergab, was mir sehr bemerkenswert erscheint.

Die Farbzerstreuung bei den Randstrahlen und den achsennächsten Strahlen lässt sich also tatsächlich vollständig aufheben. Genauer gesagt lässt sich ein Objektiv berechnen, das alle Strahlen zweier bestimmter Farben in ein und demselben Punkt vereinigt, sowohl die in einer bestimmten Entfernung von der Achse als auch die unendlich nahe an ihr auftreffenden Strahlen, wobei hier stets vorausgesetzt wird, dass sie parallel zur Achse einfallen. Dieses Objektiv erhält eine Form, die von den bisher ausschließlich verwendeten ganz abweicht. Beide Linsen werden konvex-konkav, und ihre konvexen Flächen sind dem Gegenstand zugewandt. Obwohl dabei größere Brechungen als bei anderen Anordnungen auftreten, bleibt die unvermeidliche Abweichung durch die Gestalt immer noch sehr gering. Daher werden \textit{alle} parallel zur Achse auf das Objektiv einfallenden Strahlen vollkommener vereinigt als bei jeder anderen Anordnung. Es wäre daher wohl der Mühe wert, dass geschickte Optiker diese neue Form erproben. Möglicherweise gibt es dabei derzeit noch \textit{praktische} Schwierigkeiten. Eine davon wird sein, dass die Glasstücke, aus denen die Linsen geschliffen werden sollen, dicker sein müssen. Bei der fortschreitenden Vervollkommnung der technischen Seite der Dioptrik darf man jedoch hoffen, dass sich solche Schwierigkeiten überwinden lassen. Dann ist es an der Mathematik, die ideale Form für die vollkommenste Vereinigung anzugeben.

Meine Rechnung soll lediglich als Beispiel dienen, das Gesagte zu bestätigen. Ich will die Optiker damit keineswegs an diese genauen Maße binden. Für die Glasarten, aus denen ein vollkommenes Objektiv geschliffen werden soll, müssen zunächst die Brechungs- und Zerstreuungsverhältnisse eigens und möglichst genau bestimmt und die Maße des Objektivs danach neu berechnet werden. Meiner Rechnung habe ich genau dieselben Zahlen zugrunde gelegt, mit denen
\SourceLeaf{513}\expandafter\def\csname HistoricalNoteText513\endcsname{[Dicke der ersten Linse $=200$, der zweiten $=80$ Abstand zwischen beiden Linsen $=50$. Exponenten der Brechungsverhältnisse jeweils für Kronglas und Flintglas]
\[
\left\{\begin{array}{lll}
1.525976&1.62173&\text{viol. Strahlen}\\
1.515162&1.60177&\text{mittl.}\quad\text{—}\\
1.504348&1.58181&\text{rote}\quad\text{—}
\end{array}\right\}
\]
}Professor \textsc{Bohnenberger} gerechnet hat. Auch die Dicke und den Abstand der Linsen habe ich beibehalten\HistoricalNote{513}. Da aber bei der neuen Anordnung die konvexe Fläche der Flintglaslinse eine stärkere Krümmung hat als die konkave Fläche der Kronglaslinse, können beide Linsen näher aneinander rücken (was auch in einer anderen, hier nicht weiter auszuführenden Hinsicht vorteilhaft sein wird). Ja, wenn die Optiker sonst keine Bedenken dagegen haben, kann der Zwischenraum ganz wegfallen, oder die Linsen können sich auf der Achse berühren. Das wird selbstverständlich eine gewisse Änderung der Krümmungsradien nach sich ziehen.

Ich beabsichtige nicht, den mathematischen Teil dieser Untersuchung hier auszuführen. Ich bemerke nur, dass die Aufgabe auf eine Gleichung vierten Grades mit zwei reellen Wurzeln führt, wenn man die Abweichung durch die Gestalt nach \textsc{Eulers} Art betrachtet und Dicke und Abstand der Glaslinse außer Betracht lässt. Die daraus gewonnene Näherungslösung dient als Grundlage einer indirekten Rechnung, durch die alles genau in Übereinstimmung gebracht wird. Für Mathematiker wird dieser Hinweis genügen. Die eine reelle Wurzel dieser Gleichung muss übrigens verworfen werden, weil mit ihr zu starke Krümmungen der Glasflächen verbunden sind und die unvollkommene Aufhebung der Abweichung durch die Gestalt zu deutlich spürbar würde.

Das Ergebnis meiner Rechnung lautet wie folgt.

Setzt man die Radien der Reihe nach gleich
\[
+3415{,}287;\quad -10133{,}007;\quad +4207{,}421;\quad -2807{,}320
\]
so vereinigen sich die roten und violetten Strahlen, sowohl die unendlich nahe an der Achse als auch die in der Entfernung $1083{,}687$ einfallenden, alle in einem Punkt der Achse, dessen Abstand von der letzten Fläche $=28293{,}3$ wird. Nimmt man den Abstand von der Achse, bei dem der Einfallswinkel $18^0\,30'$ beträgt, als Öffnungsradius, so beträgt der Durchmesser der Öffnung sehr nahe $\tfrac{1}{13}$ der Brennweite. Um beurtei\SourceLeaf{514}len zu können, wie groß die verbleibende Abweichung durch die Gestalt für die Strahlen zwischen dem Rand und der Achse wird, habe ich die Entfernungen, in denen sich die Strahlen bei einem Einfallswinkel von $13^0$ vereinigen, berechnet und gefunden
\[
\begin{array}{ll}
28289{,}3&\text{für die roten}\\
28290{,}0&\text{für die violetten Strahlen.}
\end{array}
\]

Ich kann nicht umhin, hier noch eine Bemerkung zu einer Äußerung von Professor \textsc{Bohnenberger} in dem erwähnten Aufsatz hinzuzufügen. Ich halte es nämlich für am vorteilhaftesten, die Abweichung durch die Gestalt für die Randstrahlen genau aufzuheben. Professor \textsc{Bohnenberger} hat dieses Verfahren auf S.\ 279, wie mir scheint, zu Unrecht getadelt. Wenn die Abweichung für die Randstrahlen genau aufgehoben sei, könne man, so sagt er, die Öffnung ohne Verlust an Deutlichkeit so weit vergrößern, bis die Abweichung wieder genauso groß werde wie die größte Abweichung der Strahlen dazwischen. Es sei daher am vorteilhaftesten, die Abweichung für Strahlen zwischen dem Rand und der Achse aufzuheben, statt für die Randstrahlen. Das wäre in der Tat richtig, wenn die verbleibende Abweichung jenseits und diesseits der Entfernung, für die sie aufgehoben ist, \textit{dasselbe Vorzeichen} hätte. Das ist jedoch nicht der Fall. Dagegen ließe sich mit einiger Plausibilität einwenden, dass es bei der Längsabweichung auf das Vorzeichen überhaupt nicht ankomme und dass positive und negative Abweichungen im Auge ein und dieselbe Unschärfe hervorriefen. Dabei würde man aber offensichtlich stillschweigend annehmen, dass das Okular stets genau so eingestellt sei, dass man das von den achsennächsten Strahlen erzeugte Bild deutlich sieht. Das lässt sich nicht zugestehen. Man kann dieses Bild immerhin das Hauptbild nennen. Es fällt mit dem von den Randstrahlen erzeugten Bild zusammen, wenn deren Abweichung aufgehoben ist. Alle übrigen Bilder liegen dann, zumindest im Allgemeinen, jenseits oder diesseits des Hauptbildes. Da man das Okular nun immer so einstellt, dass die Unschärfe möglichst klein wird, so sieht. man gerade das Hauptbild am wenigsten deutlich, und jede Vergrößerung der Öffnung vergrößert auch die Unschärfe. Eine ausführlichere Erörterung dieses Umstands würde mich hier zu weit führen.
